%%% FIELD proof-exact-feedback-cone-all-laws
Fix \(c>0\), \(P\in\Delta^{15}\), and \(\Delta\in\mathbb R\).  The declared common denominator satisfies
\[
 D_c(u)=(1+u)^2(1+cu)^2>0\qquad(u>0). \tag{1}
\]
By the chain rule, the definition of the feedback kernel, and \(\mathrm{ass{:}logit\text{-}response}\), the history-conditioned sequential law factors as
\[
 \Pr(Y_1=y,X_2=z,Y_2=w\mid X_1=x,u)
 =p_x(y,u)g_{xy}(z,u)p_z(w,u). \tag{2}
\]
Thus (2) uses only the two declared response kernels at their realized histories; it introduces no cross-history or cross-world independence.

Suppose first that \((\log c,\Delta)\in\Theta_{\mathrm w}(P)\).  By the definition of \(\Theta_{\mathrm w}\), some \((\mu,h,g)\in\mathcal M_{\mathrm w}(c)\) generates \(P\) and \(\Delta\).  Applying \(\mathrm{lem{:}response\text{-}measure\text{-}equivalence}\) and \(\mathrm{def{:}joint\text{-}response\text{-}measures}\) gives eight finite nonnegative measures \((\nu_s)_{s\in S}\), with \(\sum_s\nu_s=\mu\), such that
\[
 P(x,y,z,w)=\sum_{r:r_y=z}\int p_x(y,u)p_z(w,u)\,d\nu_{(x,r)}(u),
 \qquad
 \Delta=\sum_{s\in S}\int
 \frac{(c-1)u}{(1+u)(1+cu)}\,d\nu_s(u). \tag{3}
\]
Define \(d\rho_s=D_c^{-1}d\nu_s\), as in \(\mathrm{def{:}denominator\text{-}cleared\text{-}moments}\).  For \(0\le k\le4\), the function \(u^k/D_c(u)\) is bounded on \((0,\infty)\): it has a finite limit at zero, is continuous on compact positive intervals, and is \(O(u^{k-4})\) at infinity.  Hence all five moments \(m_{s,k}=\int u^k\,d\rho_s\) are finite.  Direct algebra from the definition of \(p\) gives
\[
 D_c(u)p_x(y,u)p_z(w,u)
 =c^{xy+zw}u^{y+w}(1+u)^{x+z}(1+cu)^{2-x-z}
 =A_{xz,yw}(u), \tag{4}
\]
and
\[
 D_c(u)\frac{(c-1)u}{(1+u)(1+cu)}
 =(c-1)u(1+u)(1+cu)=N_c(u). \tag{5}
\]
Both polynomials have degree at most four.  Substitution of (4)--(5) into (3), followed by the definition of \(L_{m_s}\), yields
\[
 P(x,y,z,w)=\sum_{r:r_y=z}L_{m_{(x,r)}}(A_{xz,yw}),
 \qquad
 \Delta=\sum_{s\in S}L_{m_s}(N_c). \tag{6}
\]
For every \(s\), the representing measure \(\rho_s\) and \(\mathrm{lem{:}open\text{-}stieltjes\text{-}quartic}\) supply \(a_s,b_s\in\mathbb R\) with \(H_s\succeq0\) and \(K_s\succeq0\).  Equations (6) and \(\mathrm{def{:}weak\text{-}lift}\) therefore give \(\mathcal F_c^{\mathrm w}(P,\Delta)\ne\varnothing\).

If the generating model is in \(\mathcal M_+(c)\), then \(\mathrm{ass{:}strict\text{-}initial}\) and \(\mathrm{ass{:}strict\text{-}feedback}\), as incorporated in \(\mathrm{def{:}strict\text{-}legal\text{-}class}\), imply for every \(s=(x,r_0,r_1)\) that
\[
 h_x(u)g_{x0}(r_0,u)g_{x1}(r_1,u)>0
 \quad\text{for }\mu\text{-almost every }u. \tag{7}
\]
By (1), each \(\rho_s\) is then mutually absolutely continuous with \(\mu\).  For \(v=(v_0,v_1,v_2)^{\mathsf T}\), let \(q_v(u)=v_0+v_1u+v_2u^2\).  The moment definition gives
\[
 v^{\mathsf T}H_sv=\int q_v(u)^2\,d\rho_s(u). \tag{8}
\]
Mutual absolute continuity makes the null space in (8) the same for all \(s\), so all \(H_s\) have the same range.  Since the null space of a sum of positive-semidefinite matrices is the intersection of their null spaces, \(H_\Sigma\) has that same range.  The equivalence in \(\mathrm{lem{:}generalized\text{-}schur\text{-}range}\) consequently supplies \(t_s\in\mathbb R\) such that \(R_s\succeq0\) for every \(s\).  By \(\mathrm{def{:}strict\text{-}lift}\), \(\mathcal F_c^+(P,\Delta)\ne\varnothing\).

Conversely, suppose \((m,a,b)\in\mathcal F_c^{\mathrm w}(P,\Delta)\).  By \(\mathrm{lem{:}open\text{-}stieltjes\text{-}quartic}\), for every \(s\) there is a finite nonnegative measure \(\rho_s\) on \((0,\infty)\) representing \(m_{s,0},\ldots,m_{s,4}\).  Set
\[
 d\nu_s(u)=D_c(u)\,d\rho_s(u). \tag{9}
\]
These measures are finite because \(D_c\) has degree four and the moments through degree four are finite.  Sum the sixteen observed-cell equalities in \(\mathrm{def{:}weak\text{-}lift}\).  For a fixed \(s=(x,r)\), the restriction \(r_y=z\) chooses exactly one \(z\) for each \(y\), and the declared logit probabilities satisfy \(\sum_y p_x(y,u)=\sum_w p_{r_y}(w,u)=1\).  Hence (4) implies that the polynomial contributed by this type after summing all paths is \(D_c\).  Since \(P\in\Delta^{15}\),
\[
 1=\sum_{x,y,z,w\in B}P(x,y,z,w)
 =\sum_{s\in S}L_{m_s}(D_c)
 =\sum_{s\in S}\nu_s((0,\infty)). \tag{10}
\]
Thus \(\sum_s\nu_s\) is a probability measure.  Equations (4)--(6) and (9) show that the path and APE integrals of this family are exactly \(P\) and \(\Delta\).  The reverse implication of \(\mathrm{lem{:}response\text{-}measure\text{-}equivalence}\) now sets \(\mu=\sum_s\nu_s\) and constructs normalized measurable kernels \(h\) and \(g\) by Radon--Nikodym derivatives.  Together with (2), this gives a member of \(\mathcal M_{\mathrm w}(c)\) generating exactly \(P\) and \(\Delta\).  Notice that (10) used only total mass one; a zero observed cell or a zero weak type measure creates no division, and probability-vector versions of the kernels may be chosen arbitrarily on the corresponding null sets.

For the strict converse, suppose \((m,a,b,t)\in\mathcal F_c^+(P,\Delta)\).  From \(R_s\succeq0\) and \(\mathrm{lem{:}generalized\text{-}schur\text{-}range}\),
\[
 \operatorname{Range}(H_\Sigma)\subseteq\operatorname{Range}(H_s)
 \qquad(s\in S). \tag{11}
\]
On the other hand, \(H_\Sigma-H_s=\sum_{j\ne s}H_j\succeq0\).  Therefore \(\ker H_\Sigma\subseteq\ker H_s\), and symmetry gives \(\operatorname{Range}(H_s)\subseteq\operatorname{Range}(H_\Sigma)\).  Hence all the ranges in (11) are equal.  If some \(H_s=0\), (11) forces \(H_\Sigma=0\); positivity of every summand then forces every \(H_j=0\), so all five moments of every type are zero.  This contradicts (10).  Thus \(m_s\ne0\) for every \(s\).

The hypotheses of \(\mathrm{lem{:}common\text{-}positive\text{-}support}\) are now met.  Choose representing measures \(\rho_s\) on one finite set \(Q\subset(0,\infty)\), every one of which gives strictly positive mass to every node of \(Q\), and define \(\nu_s\) by (9).  Positivity in (1) preserves all those strictly positive weights.  The Radon--Nikodym reconstruction at every \(q\in Q\) is explicitly
\[
 h_x(q)=
 \frac{\sum_{r\in B^2}\nu_{(x,r)}(\{q\})}
      {\sum_{j\in S}\nu_j(\{q\})}>0, \tag{12}
\]
and, for every \(x,y,z\in B\),
\[
 g_{xy}(z,q)=
 \frac{\sum_{r:r_y=z}\nu_{(x,r)}(\{q\})}
      {\sum_{r\in B^2}\nu_{(x,r)}(\{q\})}>0. \tag{13}
\]
Every denominator is positive because every type has positive weight at every common node.  Since \(\mu=\sum_s\nu_s\) is concentrated on \(Q\), (12)--(13) establish the two strict-positivity assumptions \(\mu\)-almost everywhere.  The reconstructed model therefore lies in \(\mathcal M_+(c)\) and generates exactly \(P\) and \(\Delta\).

Finally, \(m\) has \(8\cdot5=40\) coordinates and \((a,b)\) has sixteen, for fifty-six weak variables.  The weak constraints contain eight blocks \(H_s\in\mathbb S^3\) and eight blocks \(K_s\in\mathbb S^4\).  The strict lift adds eight coordinates \(t_s\) and eight blocks \(R_s\in\mathbb S^6\), for sixty-four variables.  Both implications have been proved for every probability vector \(P\), including all singular and zero-type weak regimes.  The two finite lifts therefore equal the independently model-defined fibers and are sharp, rather than outer, representations.

%%% FIELD proof-bdg-weak-fiber-translation
Let \((\mu,h,g)\in\mathcal M_{\mathrm w}(c)\) generate \(P\).  By the sequential factorization supplied by \(\mathrm{ass{:}logit\text{-}response}\), the definitions of \(h\) and \(g\), and the normalization conditions in \(\mathrm{ass{:}initial\text{-}kernel}\) and \(\mathrm{ass{:}feedback\text{-}kernel}\), summing over \((y,z,w)\) gives
\[
 q_x=\sum_{y,z,w\in B}P(x,y,z,w)=\int h_x(u)\,d\mu(u). \tag{1}
\]
The premise \(q_x>0\) therefore permits the definition, for each Borel set \(E\subset(0,\infty)\),
\[
 \pi_x(E)=q_x^{-1}\int_E h_x(u)\,d\mu(u). \tag{2}
\]
Equation (1) makes \(\pi_x\) a probability measure.  Substitution of \(h_x\,d\mu=q_x\,d\pi_x\) in the path law yields
\[
 P(x,y,z,w)
 =q_x\int p_x(y,u)g_{xy}(z,u)p_z(w,u)\,d\pi_x(u). \tag{3}
\]
If \(f_c(u)=(c-1)u/\{(1+u)(1+cu)\}\), then \(\sum_xh_x=1\), from \(\mathrm{ass{:}initial\text{-}kernel}\), also gives
\[
 \Delta=\int f_c(u)\,d\mu(u)
 =\sum_{x\in B}q_x\int f_c(u)\,d\pi_x(u). \tag{4}
\]
Equations (3)--(4) are precisely the note's measure-valued conditional reformulation of the two-period path and unconditional-APE equations.

Conversely, let \(q_0,q_1>0\) sum to one, let \(\pi_0,\pi_1\) be probability measures on \((0,\infty)\), and let \(g\) be a measurable nonnegative feedback kernel satisfying \(\sum_{z\in B}g_{xy}(z,u)=1\) for every \(x,y\in B\), almost everywhere under each \(\pi_x\).  Suppose (3)--(4) hold and set
\[
 \mu=\sum_{x\in B}q_x\pi_x. \tag{5}
\]
This is a probability measure, and \(q_x\pi_x\ll\mu\).  By \(\mathrm{lem{:}radon\text{-}nikodym\text{-}finite\text{-}family}\), measurable versions exist with
\[
 h_x=\frac{d(q_x\pi_x)}{d\mu},
 \qquad h_x\ge0,
 \qquad \sum_{x\in B}h_x=1
 \quad\mu\text{-almost everywhere}. \tag{6}
\]
A property holding \(\pi_x\)-almost everywhere holds \(\mu\)-almost everywhere on \(\{h_x>0\}\), because \(q_x\,d\pi_x=h_x\,d\mu\).  On \(\{h_x=0\}\), redefine \(g_{xy}(\cdot,u)\), if needed, to any probability vector on \(B\).  This makes the feedback normalization valid \(\mu\)-almost everywhere for every history without altering any path integral.  Equations (5)--(6), the stipulated response kernels in \(\mathrm{ass{:}logit\text{-}response}\), and the two kernel normalizations therefore define a member of \(\mathcal M_{\mathrm w}(c)\).  Since
\[
 h_x\,d\mu=q_x\,d\pi_x, \tag{7}
\]
substitution into its sequential path formula recovers (3), hence all sixteen cells, and summing (7) over \(x\) recovers (4), hence the unconditional APE.  The two constructions are inverse at the level of \((P,\Delta)\).

The transformations (2) and (5)--(7) preserve atomic, singular, and absolutely continuous components.  Accordingly, they prove equivalence only with the note's unrestricted-measure conditional reformulation.  The source-scoped description in \(\mathrm{lem{:}bdg\text{-}continuous\text{-}feedback\text{-}program}\) supplies no theorem converting arbitrary measures to a fixed-reference or Lebesgue-density representation, so no such stronger equality follows.  Finally, \(\mathrm{def{:}strict\text{-}legal\text{-}class}\) adds strict initial and feedback positivity to all conditions of \(\mathrm{def{:}weak\text{-}legal\text{-}class}\).  Hence \(\mathcal M_+(c)\subseteq\mathcal M_{\mathrm w}(c)\), and the strict fiber is an additional subclass, not a second comparator equivalence.

%%% FIELD proof-exact-feedback-cone
Fix \(c>0\), \(P\in\Delta^{15}\), \(q_0,q_1>0\), and \(\Delta\in\mathbb R\).  Applying \(\mathrm{thm{:}exact\text{-}feedback\text{-}cone\text{-}all\text{-}laws}\), whose observable-law premise is only \(P\in\Delta^{15}\), gives directly
\[
 (\log c,\Delta)\in\Theta_{\mathrm w}(P)
 \Longleftrightarrow \mathcal F_c^{\mathrm w}(P,\Delta)\ne\varnothing,
 \qquad
 (\log c,\Delta)\in\Theta_+(P)
 \Longleftrightarrow \mathcal F_c^+(P,\Delta)\ne\varnothing. \tag{1}
\]
Thus the positive-margin premise is not used to establish either finite-lift equivalence, and (1) is equality with the independently defined legal fibers rather than an outer approximation.

The premise \(q_x>0\) for both \(x\in B\) is used exactly in \(\mathrm{lem{:}bdg\text{-}weak\text{-}fiber\text{-}translation}\).  That lemma normalizes \(h_x\,d\mu\) by \(q_x\) in one direction and, in the other, mixes the two conditional laws with weights \(q_x\) and recovers \(h_x\) by Radon--Nikodym derivatives.  Its equations preserve every path probability and the unconditional APE.  Consequently, the weak model fiber in (1) equals the note's measure-valued conditional reformulation of the Bonhomme--Dano--Graham two-period equations.

The published-program scope recorded by \(\mathrm{lem{:}bdg\text{-}continuous\text{-}feedback\text{-}program}\) does not provide an absolutely continuous representation of every probability measure used here.  In fact, the conditionalization and mixing maps in the translation lemma preserve atoms and singular components.  The proved comparison therefore asserts no equality with a fixed-reference or Lebesgue-density program.

By \(\mathrm{def{:}strict\text{-}legal\text{-}class}\) and \(\mathrm{def{:}weak\text{-}legal\text{-}class}\), every strict model satisfies all weak restrictions plus strict initial-assignment and feedback positivity.  Hence \(\mathcal M_+(c)\subseteq\mathcal M_{\mathrm w}(c)\), so the strict fiber is an additional subclass rather than a second comparator match.

It remains only to verify the advertised finite sizes.  In \(\mathrm{def{:}weak\text{-}lift}\), the five moments for each of eight types contribute forty variables, and \(a\) and \(b\) contribute sixteen, totaling fifty-six; its PSD constraints are the eight \(3\times3\) blocks \(H_s\) and the eight \(4\times4\) blocks \(K_s\).  By \(\mathrm{def{:}strict\text{-}lift}\), the strict certificate adds eight variables \(t_s\) and eight \(6\times6\) blocks \(R_s\), totaling sixty-four variables.  The all-laws theorem supplies the stated forward joint-response construction and reverse legal-kernel reconstruction for arbitrary probability laws on \((0,\infty)\).  This proves every assertion of the corollary.

%%% FIELD proof-fixed-slope-endpoints
Fix \(c>0\), \(v\in\{\mathrm w,+\}\), and \(P\in\Delta^{15}\), and suppose that
\[
 I_v(P,c)=\{d\in\mathbb R:(\log c,d)\in\Theta_v(P)\} \tag{1}
\]
is nonempty.  Put \(z_{\mathrm w}=(m,a,b)\), \(z_+=(m,a,b,t)\), and
\[
 \mathcal G_v(P,c)=\{(z_v,d):z_v\in\mathcal F_c^v(P,d)\}. \tag{2}
\]
Every equality defining (2) is affine in \((z_v,d)\), by \(\mathrm{def{:}weak\text{-}lift}\) and \(\mathrm{def{:}strict\text{-}lift}\).  If \(M_0,M_1\succeq0\) and \(0\le\lambda\le1\), then for every conformable vector \(q\),
\[
 q^{\mathsf T}\{\lambda M_0+(1-\lambda)M_1\}q
 =\lambda q^{\mathsf T}M_0q+(1-\lambda)q^{\mathsf T}M_1q\ge0. \tag{3}
\]
Thus all the affine PSD constraints are convex, and \(\mathcal G_v(P,c)\) is convex.  Its linear projection onto the \(d\)-coordinate is therefore an interval.  By \(\mathrm{thm{:}exact\text{-}feedback\text{-}cone\text{-}all\text{-}laws}\), that projection is exactly \(I_v(P,c)\), including when \(P\) has zero cells.

For clarity about the strict case, if \(d\in I_+(P,c)\), a member of \(\mathcal M_+(c)\) generates \(P\).  For each path \((x,y,z,w)\in B^4\), \(\mathrm{ass{:}strict\text{-}initial}\), \(\mathrm{ass{:}strict\text{-}feedback}\), and the codomain \((0,1)\) of the declared logit kernel imply
\[
 h_x(u)p_x(y,u)g_{xy}(z,u)p_z(w,u)>0
 \quad\text{for }\mu\text{-almost every }u. \tag{4}
\]
The base-law condition makes \(\mu\) a probability measure, so integration of (4) yields \(P(x,y,z,w)>0\).  Hence a nonempty strict fiber forces all observed cells to be positive.  At a zero-cell law the all-laws strict equivalence is still exact, but both sides are empty; that boundary extension is vacuous.

By \(\mathrm{thm{:}sharp\text{-}slope\text{-}ape\text{-}envelope}\),
\[
 I_v(P,c)\subseteq
 \begin{cases}
 (0,(\sqrt c-1)/(\sqrt c+1)],&c>1,\\
 [-(1-\sqrt c)/(1+\sqrt c),0),&0<c<1,\\
 \{0\},&c=1.
 \end{cases} \tag{5}
\]
This is an outer containment for the present fixed \(P\); the cited theorem's sharpness is universal across varying admissible laws and does not imply fixed-\(P\) endpoint attainment.  In particular, the nonempty interval (1) is bounded and has finite infimum and supremum.

Expose the two optimization problems as
\[
 \underline V_v(P,c)=\inf_{d,z_v}\{d:z_v\in\mathcal F_c^v(P,d)\},
 \qquad
 \overline V_v(P,c)=\sup_{d,z_v}\{d:z_v\in\mathcal F_c^v(P,d)\}. \tag{6}
\]
Their objectives are linear, their equalities are affine, and their inequalities are affine positive-semidefinite constraints, so (6) are linear semidefinite programs.  Exactness of the all-laws cone theorem gives
\[
 \underline V_v(P,c)=\inf I_v(P,c)=L_c^v(P),
 \qquad
 \overline V_v(P,c)=\sup I_v(P,c)=U_c^v(P). \tag{7}
\]
Thus these values are sharp for the independently model-defined fiber even if no optimizer exists.

Applying the same exact equivalence at each candidate endpoint gives separately
\[
 L_c^v(P)\in I_v(P,c)
 \Longleftrightarrow
 \mathcal F_c^v(P,L_c^v(P))\ne\varnothing, \tag{8}
\]
and
\[
 U_c^v(P)\in I_v(P,c)
 \Longleftrightarrow
 \mathcal F_c^v(P,U_c^v(P))\ne\varnothing. \tag{9}
\]
These are exactly the two tests in \(\mathrm{def{:}endpoint\text{-}attainment\text{-}test}\).

Closedness of the PSD cones does not make (8)--(9) automatic, because a linear projection of a closed spectrahedron need not be closed.  Indeed,
\[
 \mathcal C=\left\{(x,y)\in\mathbb R^2:
 \begin{bmatrix}x&1\\1&y\end{bmatrix}\succeq0\right\}
\]
is closed, while the principal-minor conditions \(x\ge0\), \(y\ge0\), and \(xy\ge1\) show that its projection onto the \(x\)-axis is \((0,\infty)\).  Therefore endpoint membership must be decided by the separate feasibility tests, as asserted.

%%% FIELD proof-algebraic-decision
Fix \(v\in\{\mathrm w,+\}\), write \(z_{\mathrm w}=(m,a,b)\) and \(z_+=(m,a,b,t)\), and let
\[
 \Phi_v(P,c,d,z_v) \tag{1}
\]
denote the conjunction of \(c>0\), the sixteen observed-cell equalities, the APE equality with value \(d\), and all matrix inequalities in \(\mathcal F_c^v(P,d)\).  By \(\mathrm{lem{:}principal\text{-}minor\text{-}psd}\), every fixed-size matrix inequality in (1) is equivalent to finitely many polynomial weak inequalities.  The coefficients of \(A_{xz,yw}\) and \(N_c\) are integer polynomials in \(c\), while every entry of \(H_s\), \(K_s\), and \(R_s\) is affine in the lift variables.  Consequently, (1) is a first-order ordered-field formula whose only fixed parameters are the entries of \(P\).  No logarithm occurs: elimination uses the positive odds coordinate \(c\), after which the structural slope is \(\log c\).

For every probability vector \(P\), including a boundary law, \(\mathrm{thm{:}exact\text{-}feedback\text{-}cone\text{-}all\text{-}laws}\) gives
\[
 (\log c,d)\in\Theta_v(P)
 \Longleftrightarrow
 \exists z_v\,\Phi_v(P,c,d,z_v). \tag{2}
\]
If \(P,c,d\) are real algebraic, the right-hand side is a sentence with real algebraic coefficients.  The terminating procedure in \(\mathrm{lem{:}real\text{-}closed\text{-}field\text{-}elimination}\) returns an equivalent quantifier-free Boolean combination of polynomial sign conditions.  Exact signs of algebraic numbers decide that combination, hence decide membership.  Because principal-minor inequalities are weak, singular PSD faces are included.

Now suppose only that the entries of \(P\) are real algebraic.  The two requested coordinate projections are
\[
 \mathcal C_v(P)=\{c:\exists d,z_v\,\Phi_v(P,c,d,z_v)\},
 \qquad
 \mathcal D_v(P)=\{d:\exists c,z_v\,\Phi_v(P,c,d,z_v)\}. \tag{3}
\]
The condition \(c>0\) is already part of \(\Phi_v\).  Eliminating the quantified variables in (3) terminates by the real-closed-field lemma and produces exact semialgebraic descriptions of both projections.  The first description is for \(c\), not for the generally nonalgebraic coordinate \(\theta=\log c\).

Finally, fix algebraic \(P\) and algebraic \(c>0\), and assume the fixed-\(c\) fiber is nonempty.  By \(\mathrm{thm{:}fixed\text{-}slope\text{-}endpoints}\), both endpoints are finite.  A real number \(\ell\) is the lower endpoint exactly when
\[
\begin{aligned}
 \operatorname{Inf}_v(\ell)\;\Longleftrightarrow{}&
 \bigl[\exists d,z_v\,\Phi_v(P,c,d,z_v)\bigr]\\
 &{}\wedge\bigl[\forall d,z_v\,
   \{\Phi_v(P,c,d,z_v)\Rightarrow \ell\le d\}\bigr]\\
 &{}\wedge\bigl[\forall\varepsilon\,
   \{\varepsilon>0\Rightarrow
   \exists d,z_v\,[\Phi_v(P,c,d,z_v)\wedge d<\ell+\varepsilon]\}\bigr].
\end{aligned} \tag{4}
\]
Similarly, \(r\) is the upper endpoint exactly when
\[
\begin{aligned}
 \operatorname{Sup}_v(r)\;\Longleftrightarrow{}&
 \bigl[\exists d,z_v\,\Phi_v(P,c,d,z_v)\bigr]\\
 &{}\wedge\bigl[\forall d,z_v\,
   \{\Phi_v(P,c,d,z_v)\Rightarrow d\le r\}\bigr]\\
 &{}\wedge\bigl[\forall\varepsilon\,
   \{\varepsilon>0\Rightarrow
   \exists d,z_v\,[\Phi_v(P,c,d,z_v)\wedge r-\varepsilon<d]\}\bigr].
\end{aligned} \tag{5}
\]
Quantifier elimination applied to (4)--(5) gives quantifier-free semialgebraic descriptions, over the real algebraic coefficient field generated by \(P\) and \(c\), of the singleton sets \(\{L_c^v(P)\}\) and \(\{U_c^v(P)\}\).  To see that either singleton point is algebraic, discard identically zero polynomials from its description.  If no remaining polynomial vanished at the point, every polynomial sign in the finite description would be constant on some open interval around it.  The Boolean combination would then hold throughout that interval, contradicting singletonness.  Hence a nonzero polynomial with real algebraic coefficients vanishes at each endpoint.  The quantifier-free singleton description is also an exact terminating representation of that algebraic number, proving the asserted computation claim.

The two endpoint-inclusion flags are the separate sentences
\[
 \exists z_v\,\Phi_v(P,c,L_c^v(P),z_v),
 \qquad
 \exists z_v\,\Phi_v(P,c,U_c^v(P),z_v). \tag{6}
\]
Their equivalence to model-side inclusion follows from (2), and they coincide with \(\mathrm{def{:}endpoint\text{-}attainment\text{-}test}\).  All coefficients in (6) are algebraic by the preceding argument, so elimination and exact sign evaluation decide both flags.  The projection, membership, and endpoint conclusions have thus been proved under their respective premises, with no cellwise-positivity assumption.

%%% FIELD proof-relative-interior-consistency
Fix \(c>0\) and \(v\in\{\mathrm w,+\}\), and abbreviate
\[
 C=C_c^v,
 \qquad A=\operatorname{aff}(C).
\]
The assumption \(P\in\operatorname{ri}(C)\) in \(\mathrm{ass{:}relative\text{-}interior\text{-}law}\) implies in particular that \(C\ne\varnothing\).

First, \(C\) is convex.  Take \(Q_0,Q_1\in C\), choose \(d_j\) in the fiber over \(Q_j\), and use \(\mathrm{thm{:}exact\text{-}feedback\text{-}cone\text{-}all\text{-}laws}\) to choose \(z_j\in\mathcal F_c^v(Q_j,d_j)\).  For \(0\le\lambda\le1\), all lift equalities are affine and, for every PSD block and every vector \(a\),
\[
 a^{\mathsf T}\{\lambda M_0+(1-\lambda)M_1\}a
 =\lambda a^{\mathsf T}M_0a+(1-\lambda)a^{\mathsf T}M_1a\ge0.
\]
Therefore
\[
 \lambda z_0+(1-\lambda)z_1
 \in\mathcal F_c^v\bigl(\lambda Q_0+(1-\lambda)Q_1,
 \lambda d_0+(1-\lambda)d_1\bigr). \tag{1}
\]
The reverse implication of the all-laws theorem puts the mixed law in \(C\), proving convexity for both classes.

The lower endpoint is convex on \(C\).  Indeed, for \(\varepsilon>0\), choose feasible \(d_j<L_c^v(Q_j)+\varepsilon\).  Equation (1) implies
\[
 L_c^v(\lambda Q_0+(1-\lambda)Q_1)
 \le\lambda L_c^v(Q_0)+(1-\lambda)L_c^v(Q_1)+\varepsilon,
\]
and letting \(\varepsilon\downarrow0\) proves convexity.  Choosing instead \(d_j>U_c^v(Q_j)-\varepsilon\) proves that \(U_c^v\) is concave.  By \(\mathrm{thm{:}fixed\text{-}slope\text{-}endpoints}\), both functions are finite on \(C\).  Extend \(L_c^v\) and \(-U_c^v\) by \(+\infty\) outside \(C\); they are proper convex functions with domain \(C\).  The result \(\mathrm{lem{:}relative\text{-}interior\text{-}convex\text{-}continuity}\), applied at the assumed point \(P\in\operatorname{ri}(C)\), shows that \(L_c^v\) and \(U_c^v\) are continuous at \(P\) relative to \(A\).

The relative-interior premise also implies positive observed cells, although this is derived rather than assumed.  Let \(\mu\) be unit mass at \(u=1\), put \(h_x(1)=1/2\), and put \(g_{xy}(z,1)=1/2\) for all \(x,y,z\in B\).  These kernels satisfy both legal-class definitions, including the strict conditions, and generate
\[
 P^\star(x,y,z,w)=\frac14p_x(y,1)p_z(w,1)>0. \tag{2}
\]
Thus \(P^\star\in C\).  If some coordinate \(P_j\) were zero, a relative open ball around \(P\) in \(A\) would, for sufficiently small \(\epsilon>0\), contain \(P+\epsilon(P-P^\star)\); its \(j\)-th coordinate would be \(-\epsilon P_j^\star<0\), contradicting \(C\subseteq\Delta^{15}\).  This confirms the positive-cell consequence without replacing the stronger relative-interior premise used for continuity.

It remains to prove convergence of the argument of the endpoint maps.  Under \(\mathrm{ass{:}iid\text{-}units}\), fix one cell, let \(Z_i\) be its indicator, write \(p=\mathbb E Z_i\), and set \(W_i=Z_i-p\).  Independence and centering make every term in the fourth-moment expansion vanish unless every sample index occurs at least twice.  Since \(|W_i|\le1\),
\[
 \mathbb E\!\left[\left(\frac1n\sum_{i=1}^nW_i\right)^4\right]
 =\frac{n\mathbb E[W_1^4]+3n(n-1)\{\mathbb E[W_1^2]\}^2}{n^4}
 \le\frac3{n^2}. \tag{3}
\]
Because \(\epsilon^4\mathbf 1\{|n^{-1}\sum_iW_i|>\epsilon\}\le|n^{-1}\sum_iW_i|^4\), taking expectations in this pointwise inequality and using (3) yields, for every \(\epsilon>0\),
\[
 \Pr\!\left(\left|\frac1n\sum_{i=1}^nW_i\right|>\epsilon\right)
 \le\frac3{\epsilon^4n^2}. \tag{4}
\]
The probability of some such deviation after time \(N\) is at most \(3\epsilon^{-4}\sum_{n\ge N}n^{-2}\), which tends to zero.  Intersecting the resulting probability-one events over \(\epsilon=1/k\), \(k\in\mathbb N\), and over the sixteen cells proves
\[
 \widehat P_n\longrightarrow P\qquad\text{almost surely}. \tag{5}
\]

The affine space \(A\) is closed in the finite-dimensional cell space and contains \(P\).  Since \(\widetilde P_n\) is the Euclidean projection of \(\widehat P_n\) onto \(A\), its minimizing property gives
\[
 \|\widetilde P_n-P\|_2\le\|\widehat P_n-P\|_2. \tag{6}
\]
Thus \(\widetilde P_n\to P\) almost surely.  By \(P\in\operatorname{ri}(C)\), some relative open ball in \(A\) around \(P\) is contained in \(C\); hence \(\widetilde P_n\in C\) eventually almost surely.  Relative continuity at \(P\) now gives
\[
 L_c^v(\widetilde P_n)\longrightarrow L_c^v(P),
 \qquad
 U_c^v(\widetilde P_n)\longrightarrow U_c^v(P)
 \quad\text{almost surely}. \tag{7}
\]
The argument fixes \(c\) and uses the relative-interior ball essentially; it establishes neither continuity at attainable boundary points nor joint consistency as \(c\) varies through one.

%%% FIELD proof-c-one-sanity
At \(c=1\), \(\mathrm{thm{:}sharp\text{-}slope\text{-}ape\text{-}envelope}\) gives \(\Delta=0\) for every probability law \(\mu\).  Hence every nonempty fixed-slope APE fiber is the singleton \(\{0\}\).  We prove the rank claim directly.

Expand the declared likelihood polynomial as
\[
 A_{xz,yw}(u)=\sum_{k=0}^4\gamma_{xz,yw,k}(c)u^k.
\]
The observed-law equations in \(\mathrm{def{:}weak\text{-}lift}\) have the form \(P=C(c)m\), where the \(16\times40\) coefficient matrix is
\[
 C(c)_{(x,y,z,w),(x',r_0,r_1,k)}
 =\mathbf 1\{x=x'\}\mathbf 1\{r_y=z\}\gamma_{xz,yw,k}(c). \tag{1}
\]
Order rows lexicographically by \((x,y,z,w)\), types by \((x,r_0,r_1)\), and degrees increasingly within each type.  The first indicator in (1) makes \(C(c)\) block diagonal in \(x\).  Numbering columns in each twenty-column block from zero, select columns \((0,1,2,6,7,8,10,11)\) in the \(x=0\) block and the same columns in the \(x=1\) block.  Expansion of the definition of \(A_{xz,yw}\) gives
\[
M_0=\begin{bmatrix}
1&2c&c^2&2c&c^2&0&0&0\\
0&1&2c&1&2c&c^2&0&0\\
0&0&0&0&0&0&1&c+1\\
0&0&0&0&0&0&0&c\\
0&1&2c&0&0&0&0&1\\
0&0&1&0&0&0&0&0\\
0&0&0&1&c+1&c&0&0\\
0&0&0&0&c&c^2+c&0&0
\end{bmatrix},
\quad
M_1=\begin{bmatrix}
1&c+1&c&c+1&c&0&0&0\\
0&1&c+1&1&c+1&c&0&0\\
0&0&0&0&0&0&1&2\\
0&0&0&0&0&0&0&c\\
0&c&c^2+c&0&0&0&0&c\\
0&0&c&0&0&0&0&0\\
0&0&0&c&2c&c&0&0\\
0&0&0&0&c^2&2c^2&0&0
\end{bmatrix}. \tag{2}
\]
Successive Laplace expansion along rows with a single nonzero entry gives
\[
 \det M_0=-c^2(c-1),
 \qquad
 \det M_1=-c^6(c-1). \tag{3}
\]
The selected \(16\times16\) block-diagonal minor consequently has determinant \(c^8(c-1)^2\).  Since \(C(c)\) has sixteen rows,
\[
 \operatorname{rank}C(c)=16
 \qquad(c>0,\ c\ne1). \tag{4}
\]

At \(c=1\), the history-conditioned restriction \(\mathrm{ass{:}logit\text{-}response}\) gives \(p_0(y,u)=p_1(y,u)=:p(y,u)\) at every realized history.  The sequential factorization and \(\mathrm{ass{:}feedback\text{-}kernel}\) imply, for each \(x\in B\),
\[
\begin{aligned}
 \sum_{z\in B}P(x,1,z,0)
 &=\int h_x(u)p(1,u)
      \sum_{z\in B}g_{x1}(z,u)p(0,u)\,d\mu(u)\\
 &=\int h_x(u)p(1,u)p(0,u)\,d\mu(u)\\
 &=\int h_x(u)p(0,u)
      \sum_{z\in B}g_{x0}(z,u)p(1,u)\,d\mu(u)\\
 &=\sum_{z\in B}P(x,0,z,1).
\end{aligned} \tag{5}
\]
This uses feedback normalization separately at the two realized histories and no cross-world independence.  Summing the corresponding rows of (1) gives the same identity coefficientwise: after setting \(c=1\), both denominator-cleared sides are the same polynomial representing \(p(1,u)p(0,u)\).  The resulting two left-null vectors, one for each \(x\), have disjoint row supports and hence are independent.  Therefore \(\operatorname{rank}C(1)\le14\).

For the reverse inequality, delete in each \(x\)-block the row \((x,1,1,0)\) and select within-block columns \((0,1,2,6,7,10,11)\).  Each resulting block is
\[
 Q=\begin{bmatrix}
1&2&1&2&1&0&0\\
0&1&2&1&2&0&0\\
0&0&0&0&0&1&2\\
0&0&0&0&0&0&1\\
0&1&2&0&0&0&1\\
0&0&1&0&0&0&0\\
0&0&0&0&1&0&0
\end{bmatrix}. \tag{6}
\]
Laplace expansion gives \(\det Q=1\).  Thus the corresponding \(14\times14\) minor \(\operatorname{diag}(Q,Q)\) is nonsingular, and
\[
 \operatorname{rank}C(1)=14. \tag{7}
\]

We finally identify the resulting affine-hull change.  Put
\[
 V_c=\operatorname{Range}C(c),
 \qquad
 \ell(p)=\sum_{(x,y,z,w)\in B^4}p(x,y,z,w),
\]
and
\[
 d_x(p)=\sum_{z\in B}p(x,1,z,0)-\sum_{z\in B}p(x,0,z,1). \tag{8}
\]
Every attainable law lies in \(V_c\cap\{\ell=1\}\), and (5) imposes \(d_0=d_1=0\) at \(c=1\).

To prove that these are all affine restrictions, choose four distinct positive nodes, give each of the eight measures \(\rho_s\) strictly positive weight at every node, and multiply all weights by a common positive constant so that
\[
 \sum_{s\in S}\int D_c(u)\,d\rho_s(u)=1. \tag{9}
\]
Set \(a_s=\int u^{-1}\,d\rho_s(u)\) and \(b_s=\int u^5\,d\rho_s(u)\).  Then \(H_s\) is the Gram matrix of \((1,u,u^2)\), and \(K_s\) is the Gram matrix of \((u^{-1/2},u^{1/2},u^{3/2},u^{5/2})\).  The corresponding evaluation matrices have ranks three and four on four distinct positive nodes, so \(H_s\succ0\), \(K_s\succ0\), and \(H_\Sigma\succ0\).  Choosing
\[
 t_s>\lambda_{\max}(H_\Sigma H_s^{-1}H_\Sigma)
\]
makes \(R_s\succ0\) by direct completion of the square, equivalently by its Schur complement.  This is a strictly feasible lifted point, and its induced law \(\bar P=C(c)\bar m\) has every cell positive by (2) and positivity of every type weight.

Positive definiteness of all blocks and positivity of all cells persist under sufficiently small perturbations.  Choose a linear right inverse \(J:V_c\to\mathbb R^{40}\) of \(C(c)\).  For all sufficiently small \(v\in V_c\cap\ker\ell\), set \(m(v)=\bar m+Jv\), keeping \((a,b,t)\) fixed.  Then
\[
 C(c)m(v)=\bar P+v,
 \qquad \ell(\bar P+v)=1,
\]
and every weak and strict matrix block remains positive definite.  The APE coordinate is assigned its defining linear value \(\sum_sL_{m_s(v)}(N_c)\).  Hence \(\mathrm{thm{:}exact\text{-}feedback\text{-}cone\text{-}all\text{-}laws}\) makes all such \(\bar P+v\) attainable in both classes.  It follows that
\[
 \operatorname{aff}C_c^{v_0}=\bar P+(V_c\cap\ker\ell)
 \qquad(v_0\in\{\mathrm w,+\}). \tag{10}
\]
For \(c\ne1\), (4) gives \(V_c=\mathbb R^{16}\).  For \(c=1\), the two independent identities (5), the inclusion \(V_1\subseteq\ker d_0\cap\ker d_1\), and (7) imply \(V_1=\ker d_0\cap\ker d_1\).  Therefore
\[
 \operatorname{aff}C_c^{v_0}=
 \begin{cases}
 \{p:\ell(p)=1\},&c\ne1,\\
 \{p:\ell(p)=1,\ d_0(p)=d_1(p)=0\},&c=1.
 \end{cases} \tag{11}
\]
Thus the observable coefficient rank falls from sixteen to fourteen at \(c=1\), and the affine hull used to define \(\widetilde P_n\) gains exactly two independent restrictions there.
